\paragraph{基本概念。}\label{sub:z2-basic-notions}
本节讨论带\trInvolution{}的 \trACategory，而且取的是最简单的可能意义。
\trACategory{} $\scr{A}$ 上的一个\textbf{朴素 $\ZZ/2$-作用}（naive $\ZZ/2$-action）
是一个 \trAFunctor{} $\epsilon_{\scr{A}}\colon\scr{A}\to\scr{A}$，它的阶
$\ge 2$ 的各项为零，并且满足 $\epsilon_{\scr{A}}\circ\epsilon_{\scr{A}}
=\mathrm{Id}_{\scr{A}}$。实际上，这不过是由\trObject{}上的一个\trInvolution{}与%
\trMorphism{}上的
相应映射 $\epsilon_{\scr{A}}\colon\hom_{\scr{A}}(X_0,X_1)\to
\hom_{\scr{A}}(\epsilon_{\scr{A}}X_0,\epsilon_{\scr{A}}X_1)$ 组成的，使得
\begin{equation*}
\mu^d_{\scr{A}}\bigl(\epsilon_{\scr{A}}(a_d),\ldots,\epsilon_{\scr{A}}(a_1)\bigr)
=\epsilon_{\scr{A}}\bigl(\mu^d_{\scr{A}}(a_d,\ldots,a_1)\bigr).
\end{equation*}
一个 $\epsilon_{\scr{A}}$-\textbf{不动\trObject}（fixed object）是指满足
$\epsilon_{\scr{A}}X=X$ 的\trObject。\textbf{不动 A∞-子\trCategory}（fixed
A∞-subcategory）$\scr{A}^{\mathrm{fix}}$ 是由这些\trObject{}构成的\trFull{}
A∞-子\trCategory。\textbf{不变 A∞-子\trCategory}（invariant A∞-subcategory）
$\scr{A}^{\mathrm{inv}}$ 与 $\scr{A}^{\mathrm{fix}}$ 有相同的\trObject，但其%
\trMorphism{}现在取\textbf{不变部分}（invariant part）
$\hom_{\scr{A}^{\mathrm{inv}}}(X_0,X_1)=\hom_{\scr{A}}(X_0,X_1)^{\ZZ/2}$，%
\trComposition{}映射也相应地受限。我们需要通往不变部分的这一过渡在\trHomology{}上
有良好的性态，因此从现在起假定 $\operatorname{char}(\KK)\ne 2$。

往后我们会多次遇到下面这个初等问题的事例。令 $A=H^0(\scr{A})$，并令
$\epsilon_A$ 是由 $\epsilon_{\scr{A}}$ 诱导的 $\ZZ/2$-作用。取两个
$\epsilon_A$-不动\trObject{} $X_0$、$X_1$，它们在 $A$ 中\trIsomorphic。这个%
\trIsomorphism{}何时能
做成 $\ZZ/2$-不变的？其\trObstruction{}是类
\begin{equation}\label{eq:z2-obstruction-class}
\bigl[y=\epsilon_A(x^{-1})\cdot x\bigr]\in
\{y\in\operatorname{Aut}_A(X_0): y^{-1}=\epsilon_A(y)\}/\operatorname{Aut}_A(X_0)
\end{equation}
其中 $x$ 是从 $X_0$ 到 $X_1$ 的任意一个\trIsomorphism，而
$\operatorname{Aut}_A(X_0)=\Hom^0_A(X_0,X_0)^{\times}$ 通过\trTwisted{}共轭作用在自身
之上，也就是说，$z$ 的作用是 $y\mapsto\epsilon_A(z)^{-1}yz$。最简单的情形是
$X_0$ \textbf{自由}（free）的情形，意思是 $\Hom_A(X_0,X_0)\cong R\oplus R$ 是
同一个\trGraded{}代数 $R$ 的两个副本的\trDirectSum，而 $\epsilon_A$ 通过交换这
两个\trDirectSummand{}而作用。在这种情形下，\eqref{eq:z2-obstruction-class}中
的集合由元素对 $y=(r,r^{-1})$ 组成，而其中每一个都能通过取 $z=(r^{-1},e)$ 的%
\trTwisted{}共轭\trTrivialization。于是：

\begin{lemma}\label{thm:free-fixed-invariant-isomorphic}
设我们有两个 $\epsilon_{\scr{A}}$-不动\trObject，其一已知是自由的。那么，若二者在
$H^0(\scr{A})$ 中\trIsomorphic，则它们在 $H^0(\scr{A}^{\mathrm{inv}})$ 中也%
\trIsomorphic。
\noqed
\end{lemma}
